% capitoli/05_conclusioni.tex - Capitolo 5: Conclusioni

\chapter{Conclusioni e Sviluppi Futuri}
\label{ch:conclusioni}

\section{Sintesi del Lavoro}
Il presente lavoro di tesi, inquadrato all'interno del più ampio progetto di ricerca TEXLOS \cite{texlos2026}, ha affrontato con successo la sfida di diagnosticare ed interpretare tempestivamente i colli di bottiglia nei processi ospedalieri responsabili dell'estensione critica del Length of Stay (LOS) dei pazienti. 

Attraverso una rigorosa elaborazione metodologica, si è dimostrato il \textbf{significativo} potenziale diagnostico derivante dall'innovativo accostamento tra: 
\begin{enumerate}
    \item L'approccio \textbf{Storytelling}, capace di \textbf{rappresentare} una traccia evento clinica in linguaggio naturale semantico e \textbf{mitigare le criticità legate alla dimensionalità} dei dati.
    \item Il classificatore \textbf{\texttt{bert-base-uncased}}, che è risultato \textbf{robusto rispetto alla variabilità terminologica} locale rispetto ai cloni "ClinicalBERT", raggiungendo agilmente un solido grado di accuratezza supervisionata bilanciata (circa l'88\%), mitigando al contempo il severo sbilanciamento di classe ospedaliero in virtù della \textit{Focal Loss}.
    \item Gli esplorativi algoritmi \textbf{XAI} (Explainable AI), specificamente appoggiati su logiche assiomatiche quali gli \textbf{Integrated Gradients Adattivi}.
\end{enumerate}

Quest'ultimo punto, corroborato dallo sforzo implementativo e computazionale della ricomposizione logica pre- e post-sub-tokenizzazione di BERT in azioni nosocomiali complete, rappresenta l'apporto più vitale dell'elaborato. Si fornisce concretamente alle amministrazioni sanitarie e ai reparti un cruscotto visivo e tabulare, uno strumento d'indagine in cui punteggi d'impatto rigorosi attribuiscono responsabilità precise alle anomalie ed agli specifici atti clinici, traducendo complesse matrici di gradienti profondi in diagnostica processuale esplicita.

\section{Limiti dell'Approccio}
Pur restituendo metriche solide dinanzi a configurazioni di classe disomogenee, la pipeline qui analizzata è soggetta ad intimi e non negabili limiti concettuali ed empirici.
Il primo campanello d'avviso deriva dalla fallimentare iterazione esplorativa inerente la data augmentation affidata a LLM generativi (Ollama, \texttt{phi-4}, \texttt{gemma-3}); l'incontrollabile bias strutturale dell'artificio linguistico immesso dal Transformer ha inibito catastroficamente lo strato logico e discriminatorio dello classificatore spaziale generoso. Ciò sancisce attualmente dei ristretti gradi di dipendenza logica esclusiva su determinismi documentali originali.

In linea più profondamente filosofica e scientifica, va rigorosamente ribadito il crinale epistemologico dell'Intelligenza Artificiale attuale. L'architettura proposta, le Attention Maps e la tecnica degli Integrated Gradients restituiscono di fatto \textbf{metriche e mappe d'estrema correlazione stocastica}. Identificano lo scarto induttivo, documentano come al variare dell'evento medico coincida inevitabilmente una previsione degenziale, ed evidenziano i colli di bottiglia processuali, pur \textbf{senza stabilire una relazione di causalità diretta}. In ambiente clinico, l'IA solleva la spia semantica, esplicita che l'evento \textit{A} co-occorre con un dilatamento cronico del LOS dei pazienti, ma l'imputazione causale di fondo resta di \textbf{competenza esclusiva della valutazione clinica} dello specialista.

\section{Lezioni Apprese e Generalizzabilità del Modello}

\subsection{La Sfida della Scalabilità: Il limite operativo di BERT-Large}
Un contributo critico emerso dalla fase sperimentale riguarda l'analisi della scalabilità del modello. I test condotti con l'architettura \texttt{bert-large-uncased} hanno rivelato un plateau prestazionale precoce: l'incremento del numero di parametri non ha prodotto benefici proporzionali sulla capacità discriminativa. Questo suggerisce che per un dataset di circa 7.000 tracce, la complessità di BERT-base rappresenti il limite superiore di utilità informativa, oltre il quale si riscontra una saturazione dei pattern apprendibili a fronte di costi computazionali (VRAM) non sostenibili.

\subsection{Prospettive di Generalizzabilità Clinica}
Sebbene la sperimentazione si sia concentrata su una selezione di 20 DRG (Diagnosis Related Groups), la pipeline implementata — basata sulla traduzione deterministica dei log in linguaggio naturale — risulta intrinsecamente generalizzabile ad altri contesti clinici. Il modello ha dimostrato di poter astrarre concetti medici universali senza la necessità di pre-addestramento specifico su corpus clinici (come ClinicalBERT), rendendo la soluzione facilmente esportabile in diverse unità operative.

\section{Sviluppi Futuri}
L'approccio semantico \textit{Storytelling} + IG dischiude un fronte fecondo per i futuri accostamenti nel campo del Process Mining sanitario ed applicato. Nello specifico panorama del progetto TEXLOS \cite{texlos2026}, una ramificazione potenziale concerne non tanto il consolidamento della classificazione finale (l'end-point predittivo netto), bensì l'anticipazione online della XAI. 

Sviluppi imminenti auspicano la transizione dall'esperimento "post-mortem" (traccia completa e paziente dimesso) verso una classificazione "runtime" dinamica, estraendo Integrated Gradients parziali sui checkpoint quotidiani della degenza. Simili approcci \textit{Next-Activity-Prediction} corredati da sensibilità IG ad alta densità consentiranno all'amministrazione l'aggiustamento terapeutico predittivo e reale durante e non solamente all'esodo della crisi sanitaria, inaugurando l'avvento d'un controllo manageriale ospedaliero reattivo e semanticamente autoscopico.
